% Ad Atticum 3.9 --- headnote
% ad-atticum-03-09, 13 June 58 BC, written at Thessalonica

\textit{Cicero to Atticus, written from Thessalonica on the
Ides of June 58 BC. Quintus has left his Asian governorship
before 1 May and reached Athens on 15 May. The reckoning that
awaits him at Rome will be hard --- a returning propraetor
whose brother is now exiled and whose own conduct is open to
prosecution --- and Cicero's choice has been to send him on to
Rome rather than have him come to Thessalonica. \textsection 1
gives the reason: he could not bear to see his brother, with
that soft heart of his (\textit{mollissimo animo}), in so deep
a grief, nor to display his own ruin to him, nor to be
looked on; and besides, he writes, Quintus would not have
been able to part from him. The image is the lictors put aside,
or the embrace broken by force. The bitterness of not seeing
him was the lesser of two bitternesses.}

\textit{\textsection 2 carries the most direct accusation in
the surviving exile letters. Atticus, who has gone to Pompey
in hope, is asked sharply whether he does not yet see ``by
whose hand, by whose plotting, by whose villainy'' Cicero was
ruined --- the tricolon of \textit{quorum} is Pompey,
Hortensius, and (in some readings) Crassus, the men who let
the Clodian bill stand and whose silence Cicero now reads as
betrayal. The summary is the famous antithesis \textit{non
inimici sed invidi perdiderunt}: ``it was not enemies but
men who envied us who undid us.''}
